% GENERATED BY src/python/lmi/tables/ext_asyl_descriptives.py
% -- DO NOT EDIT
%==============================================================================%
% ext_asyl_descriptives.tex
%
% Purpose : Erweiterung, Tabelle E0 -- die analytische Stichprobe nach
%           Asylstatus sowie die Uebergangsmatrix des Status. Grundlage der
%           Praemisse, dass der Rechtsstatus den Zugang zum Kurs steuert.
% Source  : SOEPremote job 244779, src/stata/remote/ext_asylum_moderator.do
%           Rohlisting: results/transcripts/
%                       src_stata_remote_ext_asylum_moderator.do.eml
%           Transkript: results/verification/ext_asylum_moderator.md
% Needs   : booktabs, tabularx (beide bereits in bachelorarbeit.sty geladen)
% Usage   : % GENERATED BY src/python/lmi/tables/ext_asyl_descriptives.py
% -- DO NOT EDIT
%==============================================================================%
% ext_asyl_descriptives.tex
%
% Purpose : Erweiterung, Tabelle E0 -- die analytische Stichprobe nach
%           Asylstatus sowie die Uebergangsmatrix des Status. Grundlage der
%           Praemisse, dass der Rechtsstatus den Zugang zum Kurs steuert.
% Source  : SOEPremote job 244779, src/stata/remote/ext_asylum_moderator.do
%           Rohlisting: results/transcripts/
%                       src_stata_remote_ext_asylum_moderator.do.eml
%           Transkript: results/verification/ext_asylum_moderator.md
% Needs   : booktabs, tabularx (beide bereits in bachelorarbeit.sty geladen)
% Usage   : % GENERATED BY src/python/lmi/tables/ext_asyl_descriptives.py
% -- DO NOT EDIT
%==============================================================================%
% ext_asyl_descriptives.tex
%
% Purpose : Erweiterung, Tabelle E0 -- die analytische Stichprobe nach
%           Asylstatus sowie die Uebergangsmatrix des Status. Grundlage der
%           Praemisse, dass der Rechtsstatus den Zugang zum Kurs steuert.
% Source  : SOEPremote job 244779, src/stata/remote/ext_asylum_moderator.do
%           Rohlisting: results/transcripts/
%                       src_stata_remote_ext_asylum_moderator.do.eml
%           Transkript: results/verification/ext_asylum_moderator.md
% Needs   : booktabs, tabularx (beide bereits in bachelorarbeit.sty geladen)
% Usage   : % GENERATED BY src/python/lmi/tables/ext_asyl_descriptives.py
% -- DO NOT EDIT
%==============================================================================%
% ext_asyl_descriptives.tex
%
% Purpose : Erweiterung, Tabelle E0 -- die analytische Stichprobe nach
%           Asylstatus sowie die Uebergangsmatrix des Status. Grundlage der
%           Praemisse, dass der Rechtsstatus den Zugang zum Kurs steuert.
% Source  : SOEPremote job 244779, src/stata/remote/ext_asylum_moderator.do
%           Rohlisting: results/transcripts/
%                       src_stata_remote_ext_asylum_moderator.do.eml
%           Transkript: results/verification/ext_asylum_moderator.md
% Needs   : booktabs, tabularx (beide bereits in bachelorarbeit.sty geladen)
% Usage   : \input{tables/ext_asyl_descriptives}
% Author  : Emre Suemer, 2026-09-01
%==============================================================================%

\makeatletter
\@ifundefined{NC@rewrite@Z}{%
  \newcolumntype{Z}{>{\raggedright\arraybackslash\hangindent=1.9em\hangafter=1}X}%
}{}
\makeatother

\begin{table}[htbp]
  \centering
  \caption{Analytische Stichprobe nach Status des Asylantrags}
  \label{tab:ext-asyl-deskriptiv}
  \footnotesize
  \setlength{\tabcolsep}{4pt}%
  \begin{tabularx}{\textwidth}{@{}Zrrrr@{}}
    \toprule
     & Keine & & & \\
     & Entscheidung & Anerkannt & Abgelehnt & Gesamt \\
    \midrule
    Personenjahre & 2.140 (39{,}1\,\%) & 2.076 (38{,}0\,\%) & 1.251 (22{,}9\,\%) & 5.467 (100{,}0\,\%) \\
    Personen (Status bei Erstbeobachtung) & 1.422 (52{,}1\,\%) & 838 (30{,}7\,\%) & 470 (17{,}2\,\%) & 2.730 (100{,}0\,\%) \\
    Vertretene Zuweisungskreise & 239 & 157 & 179 & -- \\
    \midrule
    \multicolumn{5}{@{}l}{\textit{Behandlung und Aufenthaltsdauer,
      Mittelwert (SD)}} \\
    Kreisangebot an Integrationskursen, std. & 0{,}00 (0{,}99) & 0{,}01 (1{,}03) & $-$0{,}02 (0{,}97) & 0{,}00 (1{,}00) \\
    Aufenthaltsdauer, in Jahren & 2{,}12 (1{,}20) & 2{,}50 (1{,}00) & 2{,}80 (1{,}10) & 2{,}42 (1{,}13) \\
    \midrule
    \multicolumn{5}{@{}l}{\textit{Ergebnis und Mechanismen, Mittelwerte}} \\
    Erwerbst\"atig & 14{,}2\,\% & 18{,}5\,\% & 22{,}1\,\% & 17{,}6\,\% \\
    \quad \textit{N} & 2.140 & 2.076 & 1.251 & 5.467 \\
    Deutschkenntnisse (Skala) & 4{,}95 & 5{,}89 & 5{,}33 & 5{,}40 \\
    \quad \textit{N} & 2.136 & 2.076 & 1.251 & 5.463 \\
    Integrationskurs besucht & 34{,}1\,\% & 62{,}3\,\% & 40{,}7\,\% & 46{,}3\,\% \\
    \quad \textit{N} & 2.123 & 2.065 & 1.240 & 5.428 \\
    Integrationskurs abgeschlossen & 19{,}4\,\% & 38{,}4\,\% & 28{,}2\,\% & 28{,}5\,\% \\
    \quad \textit{N} & 2.062 & 1.949 & 1.186 & 5.197 \\
    Sprachzertifikat erworben & 30{,}0\,\% & 52{,}6\,\% & 48{,}8\,\% & 42{,}9\,\% \\
    \quad \textit{N} & 2.132 & 2.075 & 1.248 & 5.455 \\
    Kontakth\"aufigkeit zu Deutschen, std. & $-$0{,}01 & $-$0{,}02 & $-$0{,}01 & $-$0{,}01 \\
    \quad \textit{N} & 909 & 1.739 & 1.062 & 3.710 \\
    \bottomrule
  \end{tabularx}

  \vspace{0.75ex}
  \parbox{\textwidth}{\footnotesize
    \textit{Anmerkungen:} Eigene Berechnung, IAB-BAMF-SOEP-Befragung von
    Gefl\"uchteten 2016--2019, SOEP-Core v36 (Remote Edition), analytische
    Stichprobe der Replikation (5.467 Personenjahre, 2.730 Personen),
    ungewichtet und ohne Imputation. Der Status ist im Erhebungsjahr gemessen;
    die Personenzeile z\"ahlt den Status bei der Erstbeobachtung.
    Bin\"are Variablen sind in Prozent ausgewiesen. }
	
	%Die Kreiszeile summiert 
	%sich nicht \"uber die Spalten: ein Kreis mit Gefl\"uchteten in mehreren
	%Statusgruppen wird je Gruppe einmal gez\"ahlt. Die Fallzahlen variieren
	%zwischen den Zeilen, da die Mechanismen unterschiedlich vollst\"andig
	%beobachtet sind (Spannweite 909 bis 2.140 Personenjahre);
	%\texttt{asyl\_req\_stat} selbst ist in allen 5.467 Personenjahren
	%beobachtet und wird deshalb nirgends imputiert.
	%\\[0.4ex]
	%\textbf{Zwei Lesehinweise.} Erstens ist die Zeile \glqq Integrationskurs
	%besucht\grqq{} die Pr\"amisse der Erweiterung: 28{,}2~Prozentpunkte
	%Unterschied zwischen Anerkannten und Antragstellenden ohne Entscheidung --
	%die gr\"o\ss{}te gemessene Differenz und genau das Muster, das
	%\S~44~AufenthG erwarten l\"asst. Es handelt sich um einen
	%\emph{Haupteffekt}, nicht um die getestete Interaktion.
	%Zweitens ist die ann\"ahernd gleiche Gr\"o\ss{}e der drei Gruppen ein
	%Artefakt der Stichprobenauswahl: \texttt{clean\_sample.do:221} beh\"alt nur
	%Beobachtungen unter Wohnsitzauflage, und alle nicht mehr gebundenen
	%Beobachtungen sind anerkannte. \glqq Anerkannt\grqq{} bedeutet hier also
	%\emph{anerkannt und weiterhin gebunden}; die Verteilung ist nicht
	%repr\"asentativ f\"ur die Gefl\"uchtetenpopulation. Abgelehnte sind zudem
	%am l\"angsten in Deutschland (2{,}80 gegen\"uber 2{,}12 Jahren), die
	%Rangfolge ist daher keine reine Rangfolge des Rechtszugangs.
    
\end{table}





\begin{table}[htbp]
  \centering
  \caption{Status bei Erstbeobachtung nach Status im Erhebungsjahr}
  \label{tab:ext-asyl-uebergang}
  \footnotesize
  \setlength{\tabcolsep}{4pt}%
  \begin{tabularx}{\textwidth}{@{}Zrrrr@{}}
    \toprule
     & \multicolumn{3}{c}{Status im Erhebungsjahr} & \\
    \cmidrule(lr){2-4}
    Status bei Erstbeobachtung & Keine Entsch. & Anerkannt & Abgelehnt
     & Gesamt \\
    \midrule
    Keine Entscheidung & 2.140 (74{,}4\,\%) & 431 (15{,}0\,\%) & 307 (10{,}7\,\%) & 2.878 \\
    Anerkannt & 0 (0{,}0\,\%) & 1.645 (100{,}0\,\%) & 0 (0{,}0\,\%) & 1.645 \\
    Abgelehnt & 0 (0{,}0\,\%) & 0 (0{,}0\,\%) & 944 (100{,}0\,\%) & 944 \\
    \midrule
    Gesamt & 2.140 & 2.076 & 1.251 & 5.467 \\
    \bottomrule
  \end{tabularx}

  \vspace{0.75ex}
  \parbox{\textwidth}{\footnotesize
    \textit{Anmerkungen:} Personenjahre, Zeilenprozente in Klammern.
    Quelle wie Tabelle~\ref{tab:ext-asyl-deskriptiv}.
    Statuswechsel sind selten und strikt absorbierend: nur 369 von 2.730
    Personen (13{,}5\,\%) wechseln \"uberhaupt, und Wechsel gehen
    ausschlie\ss{}lich von \glqq keine Entscheidung\grqq{} aus -- einmal
    anerkannt oder abgelehnt, bleibt der Status bestehen.

	%    \\[0.4ex]
	%\textbf{Warum das f\"ur die Spezifikation z\"ahlt:} Asylentscheidungen
	%ergehen \emph{nach} der Zuweisung auf einen Kreis, die die Behandlung
	%definiert. Der im Erhebungsjahr gemessene Status ist damit f\"ur genau die
	%wechselnden Personen ein Nachbehandlungsmerkmal. Die Sch\"atzung wird
	%deshalb zweifach berichtet -- mit Status im Erhebungsjahr (vergleichbar zur
	%Vorlage) und mit Status bei Erstbeobachtung (konzeptionell sauberer); siehe
	%Tabelle~\ref{tab:ext-asyl-erwerb}.
}
\end{table}

% Author  : Emre Suemer, 2026-09-01
%==============================================================================%

\makeatletter
\@ifundefined{NC@rewrite@Z}{%
  \newcolumntype{Z}{>{\raggedright\arraybackslash\hangindent=1.9em\hangafter=1}X}%
}{}
\makeatother

\begin{table}[htbp]
  \centering
  \caption{Analytische Stichprobe nach Status des Asylantrags}
  \label{tab:ext-asyl-deskriptiv}
  \footnotesize
  \setlength{\tabcolsep}{4pt}%
  \begin{tabularx}{\textwidth}{@{}Zrrrr@{}}
    \toprule
     & Keine & & & \\
     & Entscheidung & Anerkannt & Abgelehnt & Gesamt \\
    \midrule
    Personenjahre & 2.140 (39{,}1\,\%) & 2.076 (38{,}0\,\%) & 1.251 (22{,}9\,\%) & 5.467 (100{,}0\,\%) \\
    Personen (Status bei Erstbeobachtung) & 1.422 (52{,}1\,\%) & 838 (30{,}7\,\%) & 470 (17{,}2\,\%) & 2.730 (100{,}0\,\%) \\
    Vertretene Zuweisungskreise & 239 & 157 & 179 & -- \\
    \midrule
    \multicolumn{5}{@{}l}{\textit{Behandlung und Aufenthaltsdauer,
      Mittelwert (SD)}} \\
    Kreisangebot an Integrationskursen, std. & 0{,}00 (0{,}99) & 0{,}01 (1{,}03) & $-$0{,}02 (0{,}97) & 0{,}00 (1{,}00) \\
    Aufenthaltsdauer, in Jahren & 2{,}12 (1{,}20) & 2{,}50 (1{,}00) & 2{,}80 (1{,}10) & 2{,}42 (1{,}13) \\
    \midrule
    \multicolumn{5}{@{}l}{\textit{Ergebnis und Mechanismen, Mittelwerte}} \\
    Erwerbst\"atig & 14{,}2\,\% & 18{,}5\,\% & 22{,}1\,\% & 17{,}6\,\% \\
    \quad \textit{N} & 2.140 & 2.076 & 1.251 & 5.467 \\
    Deutschkenntnisse (Skala) & 4{,}95 & 5{,}89 & 5{,}33 & 5{,}40 \\
    \quad \textit{N} & 2.136 & 2.076 & 1.251 & 5.463 \\
    Integrationskurs besucht & 34{,}1\,\% & 62{,}3\,\% & 40{,}7\,\% & 46{,}3\,\% \\
    \quad \textit{N} & 2.123 & 2.065 & 1.240 & 5.428 \\
    Integrationskurs abgeschlossen & 19{,}4\,\% & 38{,}4\,\% & 28{,}2\,\% & 28{,}5\,\% \\
    \quad \textit{N} & 2.062 & 1.949 & 1.186 & 5.197 \\
    Sprachzertifikat erworben & 30{,}0\,\% & 52{,}6\,\% & 48{,}8\,\% & 42{,}9\,\% \\
    \quad \textit{N} & 2.132 & 2.075 & 1.248 & 5.455 \\
    Kontakth\"aufigkeit zu Deutschen, std. & $-$0{,}01 & $-$0{,}02 & $-$0{,}01 & $-$0{,}01 \\
    \quad \textit{N} & 909 & 1.739 & 1.062 & 3.710 \\
    \bottomrule
  \end{tabularx}

  \vspace{0.75ex}
  \parbox{\textwidth}{\footnotesize
    \textit{Anmerkungen:} Eigene Berechnung, IAB-BAMF-SOEP-Befragung von
    Gefl\"uchteten 2016--2019, SOEP-Core v36 (Remote Edition), analytische
    Stichprobe der Replikation (5.467 Personenjahre, 2.730 Personen),
    ungewichtet und ohne Imputation. Der Status ist im Erhebungsjahr gemessen;
    die Personenzeile z\"ahlt den Status bei der Erstbeobachtung.
    Bin\"are Variablen sind in Prozent ausgewiesen. }
	
	%Die Kreiszeile summiert 
	%sich nicht \"uber die Spalten: ein Kreis mit Gefl\"uchteten in mehreren
	%Statusgruppen wird je Gruppe einmal gez\"ahlt. Die Fallzahlen variieren
	%zwischen den Zeilen, da die Mechanismen unterschiedlich vollst\"andig
	%beobachtet sind (Spannweite 909 bis 2.140 Personenjahre);
	%\texttt{asyl\_req\_stat} selbst ist in allen 5.467 Personenjahren
	%beobachtet und wird deshalb nirgends imputiert.
	%\\[0.4ex]
	%\textbf{Zwei Lesehinweise.} Erstens ist die Zeile \glqq Integrationskurs
	%besucht\grqq{} die Pr\"amisse der Erweiterung: 28{,}2~Prozentpunkte
	%Unterschied zwischen Anerkannten und Antragstellenden ohne Entscheidung --
	%die gr\"o\ss{}te gemessene Differenz und genau das Muster, das
	%\S~44~AufenthG erwarten l\"asst. Es handelt sich um einen
	%\emph{Haupteffekt}, nicht um die getestete Interaktion.
	%Zweitens ist die ann\"ahernd gleiche Gr\"o\ss{}e der drei Gruppen ein
	%Artefakt der Stichprobenauswahl: \texttt{clean\_sample.do:221} beh\"alt nur
	%Beobachtungen unter Wohnsitzauflage, und alle nicht mehr gebundenen
	%Beobachtungen sind anerkannte. \glqq Anerkannt\grqq{} bedeutet hier also
	%\emph{anerkannt und weiterhin gebunden}; die Verteilung ist nicht
	%repr\"asentativ f\"ur die Gefl\"uchtetenpopulation. Abgelehnte sind zudem
	%am l\"angsten in Deutschland (2{,}80 gegen\"uber 2{,}12 Jahren), die
	%Rangfolge ist daher keine reine Rangfolge des Rechtszugangs.
    
\end{table}





\begin{table}[htbp]
  \centering
  \caption{Status bei Erstbeobachtung nach Status im Erhebungsjahr}
  \label{tab:ext-asyl-uebergang}
  \footnotesize
  \setlength{\tabcolsep}{4pt}%
  \begin{tabularx}{\textwidth}{@{}Zrrrr@{}}
    \toprule
     & \multicolumn{3}{c}{Status im Erhebungsjahr} & \\
    \cmidrule(lr){2-4}
    Status bei Erstbeobachtung & Keine Entsch. & Anerkannt & Abgelehnt
     & Gesamt \\
    \midrule
    Keine Entscheidung & 2.140 (74{,}4\,\%) & 431 (15{,}0\,\%) & 307 (10{,}7\,\%) & 2.878 \\
    Anerkannt & 0 (0{,}0\,\%) & 1.645 (100{,}0\,\%) & 0 (0{,}0\,\%) & 1.645 \\
    Abgelehnt & 0 (0{,}0\,\%) & 0 (0{,}0\,\%) & 944 (100{,}0\,\%) & 944 \\
    \midrule
    Gesamt & 2.140 & 2.076 & 1.251 & 5.467 \\
    \bottomrule
  \end{tabularx}

  \vspace{0.75ex}
  \parbox{\textwidth}{\footnotesize
    \textit{Anmerkungen:} Personenjahre, Zeilenprozente in Klammern.
    Quelle wie Tabelle~\ref{tab:ext-asyl-deskriptiv}.
    Statuswechsel sind selten und strikt absorbierend: nur 369 von 2.730
    Personen (13{,}5\,\%) wechseln \"uberhaupt, und Wechsel gehen
    ausschlie\ss{}lich von \glqq keine Entscheidung\grqq{} aus -- einmal
    anerkannt oder abgelehnt, bleibt der Status bestehen.

	%    \\[0.4ex]
	%\textbf{Warum das f\"ur die Spezifikation z\"ahlt:} Asylentscheidungen
	%ergehen \emph{nach} der Zuweisung auf einen Kreis, die die Behandlung
	%definiert. Der im Erhebungsjahr gemessene Status ist damit f\"ur genau die
	%wechselnden Personen ein Nachbehandlungsmerkmal. Die Sch\"atzung wird
	%deshalb zweifach berichtet -- mit Status im Erhebungsjahr (vergleichbar zur
	%Vorlage) und mit Status bei Erstbeobachtung (konzeptionell sauberer); siehe
	%Tabelle~\ref{tab:ext-asyl-erwerb}.
}
\end{table}

% Author  : Emre Suemer, 2026-09-01
%==============================================================================%

\makeatletter
\@ifundefined{NC@rewrite@Z}{%
  \newcolumntype{Z}{>{\raggedright\arraybackslash\hangindent=1.9em\hangafter=1}X}%
}{}
\makeatother

\begin{table}[htbp]
  \centering
  \caption{Analytische Stichprobe nach Status des Asylantrags}
  \label{tab:ext-asyl-deskriptiv}
  \footnotesize
  \setlength{\tabcolsep}{4pt}%
  \begin{tabularx}{\textwidth}{@{}Zrrrr@{}}
    \toprule
     & Keine & & & \\
     & Entscheidung & Anerkannt & Abgelehnt & Gesamt \\
    \midrule
    Personenjahre & 2.140 (39{,}1\,\%) & 2.076 (38{,}0\,\%) & 1.251 (22{,}9\,\%) & 5.467 (100{,}0\,\%) \\
    Personen (Status bei Erstbeobachtung) & 1.422 (52{,}1\,\%) & 838 (30{,}7\,\%) & 470 (17{,}2\,\%) & 2.730 (100{,}0\,\%) \\
    Vertretene Zuweisungskreise & 239 & 157 & 179 & -- \\
    \midrule
    \multicolumn{5}{@{}l}{\textit{Behandlung und Aufenthaltsdauer,
      Mittelwert (SD)}} \\
    Kreisangebot an Integrationskursen, std. & 0{,}00 (0{,}99) & 0{,}01 (1{,}03) & $-$0{,}02 (0{,}97) & 0{,}00 (1{,}00) \\
    Aufenthaltsdauer, in Jahren & 2{,}12 (1{,}20) & 2{,}50 (1{,}00) & 2{,}80 (1{,}10) & 2{,}42 (1{,}13) \\
    \midrule
    \multicolumn{5}{@{}l}{\textit{Ergebnis und Mechanismen, Mittelwerte}} \\
    Erwerbst\"atig & 14{,}2\,\% & 18{,}5\,\% & 22{,}1\,\% & 17{,}6\,\% \\
    \quad \textit{N} & 2.140 & 2.076 & 1.251 & 5.467 \\
    Deutschkenntnisse (Skala) & 4{,}95 & 5{,}89 & 5{,}33 & 5{,}40 \\
    \quad \textit{N} & 2.136 & 2.076 & 1.251 & 5.463 \\
    Integrationskurs besucht & 34{,}1\,\% & 62{,}3\,\% & 40{,}7\,\% & 46{,}3\,\% \\
    \quad \textit{N} & 2.123 & 2.065 & 1.240 & 5.428 \\
    Integrationskurs abgeschlossen & 19{,}4\,\% & 38{,}4\,\% & 28{,}2\,\% & 28{,}5\,\% \\
    \quad \textit{N} & 2.062 & 1.949 & 1.186 & 5.197 \\
    Sprachzertifikat erworben & 30{,}0\,\% & 52{,}6\,\% & 48{,}8\,\% & 42{,}9\,\% \\
    \quad \textit{N} & 2.132 & 2.075 & 1.248 & 5.455 \\
    Kontakth\"aufigkeit zu Deutschen, std. & $-$0{,}01 & $-$0{,}02 & $-$0{,}01 & $-$0{,}01 \\
    \quad \textit{N} & 909 & 1.739 & 1.062 & 3.710 \\
    \bottomrule
  \end{tabularx}

  \vspace{0.75ex}
  \parbox{\textwidth}{\footnotesize
    \textit{Anmerkungen:} Eigene Berechnung, IAB-BAMF-SOEP-Befragung von
    Gefl\"uchteten 2016--2019, SOEP-Core v36 (Remote Edition), analytische
    Stichprobe der Replikation (5.467 Personenjahre, 2.730 Personen),
    ungewichtet und ohne Imputation. Der Status ist im Erhebungsjahr gemessen;
    die Personenzeile z\"ahlt den Status bei der Erstbeobachtung.
    Bin\"are Variablen sind in Prozent ausgewiesen. }
	
	%Die Kreiszeile summiert 
	%sich nicht \"uber die Spalten: ein Kreis mit Gefl\"uchteten in mehreren
	%Statusgruppen wird je Gruppe einmal gez\"ahlt. Die Fallzahlen variieren
	%zwischen den Zeilen, da die Mechanismen unterschiedlich vollst\"andig
	%beobachtet sind (Spannweite 909 bis 2.140 Personenjahre);
	%\texttt{asyl\_req\_stat} selbst ist in allen 5.467 Personenjahren
	%beobachtet und wird deshalb nirgends imputiert.
	%\\[0.4ex]
	%\textbf{Zwei Lesehinweise.} Erstens ist die Zeile \glqq Integrationskurs
	%besucht\grqq{} die Pr\"amisse der Erweiterung: 28{,}2~Prozentpunkte
	%Unterschied zwischen Anerkannten und Antragstellenden ohne Entscheidung --
	%die gr\"o\ss{}te gemessene Differenz und genau das Muster, das
	%\S~44~AufenthG erwarten l\"asst. Es handelt sich um einen
	%\emph{Haupteffekt}, nicht um die getestete Interaktion.
	%Zweitens ist die ann\"ahernd gleiche Gr\"o\ss{}e der drei Gruppen ein
	%Artefakt der Stichprobenauswahl: \texttt{clean\_sample.do:221} beh\"alt nur
	%Beobachtungen unter Wohnsitzauflage, und alle nicht mehr gebundenen
	%Beobachtungen sind anerkannte. \glqq Anerkannt\grqq{} bedeutet hier also
	%\emph{anerkannt und weiterhin gebunden}; die Verteilung ist nicht
	%repr\"asentativ f\"ur die Gefl\"uchtetenpopulation. Abgelehnte sind zudem
	%am l\"angsten in Deutschland (2{,}80 gegen\"uber 2{,}12 Jahren), die
	%Rangfolge ist daher keine reine Rangfolge des Rechtszugangs.
    
\end{table}





\begin{table}[htbp]
  \centering
  \caption{Status bei Erstbeobachtung nach Status im Erhebungsjahr}
  \label{tab:ext-asyl-uebergang}
  \footnotesize
  \setlength{\tabcolsep}{4pt}%
  \begin{tabularx}{\textwidth}{@{}Zrrrr@{}}
    \toprule
     & \multicolumn{3}{c}{Status im Erhebungsjahr} & \\
    \cmidrule(lr){2-4}
    Status bei Erstbeobachtung & Keine Entsch. & Anerkannt & Abgelehnt
     & Gesamt \\
    \midrule
    Keine Entscheidung & 2.140 (74{,}4\,\%) & 431 (15{,}0\,\%) & 307 (10{,}7\,\%) & 2.878 \\
    Anerkannt & 0 (0{,}0\,\%) & 1.645 (100{,}0\,\%) & 0 (0{,}0\,\%) & 1.645 \\
    Abgelehnt & 0 (0{,}0\,\%) & 0 (0{,}0\,\%) & 944 (100{,}0\,\%) & 944 \\
    \midrule
    Gesamt & 2.140 & 2.076 & 1.251 & 5.467 \\
    \bottomrule
  \end{tabularx}

  \vspace{0.75ex}
  \parbox{\textwidth}{\footnotesize
    \textit{Anmerkungen:} Personenjahre, Zeilenprozente in Klammern.
    Quelle wie Tabelle~\ref{tab:ext-asyl-deskriptiv}.
    Statuswechsel sind selten und strikt absorbierend: nur 369 von 2.730
    Personen (13{,}5\,\%) wechseln \"uberhaupt, und Wechsel gehen
    ausschlie\ss{}lich von \glqq keine Entscheidung\grqq{} aus -- einmal
    anerkannt oder abgelehnt, bleibt der Status bestehen.

	%    \\[0.4ex]
	%\textbf{Warum das f\"ur die Spezifikation z\"ahlt:} Asylentscheidungen
	%ergehen \emph{nach} der Zuweisung auf einen Kreis, die die Behandlung
	%definiert. Der im Erhebungsjahr gemessene Status ist damit f\"ur genau die
	%wechselnden Personen ein Nachbehandlungsmerkmal. Die Sch\"atzung wird
	%deshalb zweifach berichtet -- mit Status im Erhebungsjahr (vergleichbar zur
	%Vorlage) und mit Status bei Erstbeobachtung (konzeptionell sauberer); siehe
	%Tabelle~\ref{tab:ext-asyl-erwerb}.
}
\end{table}

% Author  : Emre Suemer, 2026-09-01
%==============================================================================%

\makeatletter
\@ifundefined{NC@rewrite@Z}{%
  \newcolumntype{Z}{>{\raggedright\arraybackslash\hangindent=1.9em\hangafter=1}X}%
}{}
\makeatother

\begin{table}[htbp]
  \centering
  \caption{Analytische Stichprobe nach Status des Asylantrags}
  \label{tab:ext-asyl-deskriptiv}
  \footnotesize
  \setlength{\tabcolsep}{4pt}%
  \begin{tabularx}{\textwidth}{@{}Zrrrr@{}}
    \toprule
     & Keine & & & \\
     & Entscheidung & Anerkannt & Abgelehnt & Gesamt \\
    \midrule
    Personenjahre & 2.140 (39{,}1\,\%) & 2.076 (38{,}0\,\%) & 1.251 (22{,}9\,\%) & 5.467 (100{,}0\,\%) \\
    Personen (Status bei Erstbeobachtung) & 1.422 (52{,}1\,\%) & 838 (30{,}7\,\%) & 470 (17{,}2\,\%) & 2.730 (100{,}0\,\%) \\
    Vertretene Zuweisungskreise & 239 & 157 & 179 & -- \\
    \midrule
    \multicolumn{5}{@{}l}{\textit{Behandlung und Aufenthaltsdauer,
      Mittelwert (SD)}} \\
    Kreisangebot an Integrationskursen, std. & 0{,}00 (0{,}99) & 0{,}01 (1{,}03) & $-$0{,}02 (0{,}97) & 0{,}00 (1{,}00) \\
    Aufenthaltsdauer, in Jahren & 2{,}12 (1{,}20) & 2{,}50 (1{,}00) & 2{,}80 (1{,}10) & 2{,}42 (1{,}13) \\
    \midrule
    \multicolumn{5}{@{}l}{\textit{Ergebnis und Mechanismen, Mittelwerte}} \\
    Erwerbst\"atig & 14{,}2\,\% & 18{,}5\,\% & 22{,}1\,\% & 17{,}6\,\% \\
    \quad \textit{N} & 2.140 & 2.076 & 1.251 & 5.467 \\
    Deutschkenntnisse (Skala) & 4{,}95 & 5{,}89 & 5{,}33 & 5{,}40 \\
    \quad \textit{N} & 2.136 & 2.076 & 1.251 & 5.463 \\
    Integrationskurs besucht & 34{,}1\,\% & 62{,}3\,\% & 40{,}7\,\% & 46{,}3\,\% \\
    \quad \textit{N} & 2.123 & 2.065 & 1.240 & 5.428 \\
    Integrationskurs abgeschlossen & 19{,}4\,\% & 38{,}4\,\% & 28{,}2\,\% & 28{,}5\,\% \\
    \quad \textit{N} & 2.062 & 1.949 & 1.186 & 5.197 \\
    Sprachzertifikat erworben & 30{,}0\,\% & 52{,}6\,\% & 48{,}8\,\% & 42{,}9\,\% \\
    \quad \textit{N} & 2.132 & 2.075 & 1.248 & 5.455 \\
    Kontakth\"aufigkeit zu Deutschen, std. & $-$0{,}01 & $-$0{,}02 & $-$0{,}01 & $-$0{,}01 \\
    \quad \textit{N} & 909 & 1.739 & 1.062 & 3.710 \\
    \bottomrule
  \end{tabularx}

  \vspace{0.75ex}
  \parbox{\textwidth}{\footnotesize
    \textit{Anmerkungen:} Eigene Berechnung, IAB-BAMF-SOEP-Befragung von
    Gefl\"uchteten 2016--2019, SOEP-Core v36 (Remote Edition), analytische
    Stichprobe der Replikation (5.467 Personenjahre, 2.730 Personen),
    ungewichtet und ohne Imputation. Der Status ist im Erhebungsjahr gemessen;
    die Personenzeile z\"ahlt den Status bei der Erstbeobachtung.
    Bin\"are Variablen sind in Prozent ausgewiesen. }
	
	%Die Kreiszeile summiert 
	%sich nicht \"uber die Spalten: ein Kreis mit Gefl\"uchteten in mehreren
	%Statusgruppen wird je Gruppe einmal gez\"ahlt. Die Fallzahlen variieren
	%zwischen den Zeilen, da die Mechanismen unterschiedlich vollst\"andig
	%beobachtet sind (Spannweite 909 bis 2.140 Personenjahre);
	%\texttt{asyl\_req\_stat} selbst ist in allen 5.467 Personenjahren
	%beobachtet und wird deshalb nirgends imputiert.
	%\\[0.4ex]
	%\textbf{Zwei Lesehinweise.} Erstens ist die Zeile \glqq Integrationskurs
	%besucht\grqq{} die Pr\"amisse der Erweiterung: 28{,}2~Prozentpunkte
	%Unterschied zwischen Anerkannten und Antragstellenden ohne Entscheidung --
	%die gr\"o\ss{}te gemessene Differenz und genau das Muster, das
	%\S~44~AufenthG erwarten l\"asst. Es handelt sich um einen
	%\emph{Haupteffekt}, nicht um die getestete Interaktion.
	%Zweitens ist die ann\"ahernd gleiche Gr\"o\ss{}e der drei Gruppen ein
	%Artefakt der Stichprobenauswahl: \texttt{clean\_sample.do:221} beh\"alt nur
	%Beobachtungen unter Wohnsitzauflage, und alle nicht mehr gebundenen
	%Beobachtungen sind anerkannte. \glqq Anerkannt\grqq{} bedeutet hier also
	%\emph{anerkannt und weiterhin gebunden}; die Verteilung ist nicht
	%repr\"asentativ f\"ur die Gefl\"uchtetenpopulation. Abgelehnte sind zudem
	%am l\"angsten in Deutschland (2{,}80 gegen\"uber 2{,}12 Jahren), die
	%Rangfolge ist daher keine reine Rangfolge des Rechtszugangs.
    
\end{table}





\begin{table}[htbp]
  \centering
  \caption{Status bei Erstbeobachtung nach Status im Erhebungsjahr}
  \label{tab:ext-asyl-uebergang}
  \footnotesize
  \setlength{\tabcolsep}{4pt}%
  \begin{tabularx}{\textwidth}{@{}Zrrrr@{}}
    \toprule
     & \multicolumn{3}{c}{Status im Erhebungsjahr} & \\
    \cmidrule(lr){2-4}
    Status bei Erstbeobachtung & Keine Entsch. & Anerkannt & Abgelehnt
     & Gesamt \\
    \midrule
    Keine Entscheidung & 2.140 (74{,}4\,\%) & 431 (15{,}0\,\%) & 307 (10{,}7\,\%) & 2.878 \\
    Anerkannt & 0 (0{,}0\,\%) & 1.645 (100{,}0\,\%) & 0 (0{,}0\,\%) & 1.645 \\
    Abgelehnt & 0 (0{,}0\,\%) & 0 (0{,}0\,\%) & 944 (100{,}0\,\%) & 944 \\
    \midrule
    Gesamt & 2.140 & 2.076 & 1.251 & 5.467 \\
    \bottomrule
  \end{tabularx}

  \vspace{0.75ex}
  \parbox{\textwidth}{\footnotesize
    \textit{Anmerkungen:} Personenjahre, Zeilenprozente in Klammern.
    Quelle wie Tabelle~\ref{tab:ext-asyl-deskriptiv}.
    Statuswechsel sind selten und strikt absorbierend: nur 369 von 2.730
    Personen (13{,}5\,\%) wechseln \"uberhaupt, und Wechsel gehen
    ausschlie\ss{}lich von \glqq keine Entscheidung\grqq{} aus -- einmal
    anerkannt oder abgelehnt, bleibt der Status bestehen.

	%    \\[0.4ex]
	%\textbf{Warum das f\"ur die Spezifikation z\"ahlt:} Asylentscheidungen
	%ergehen \emph{nach} der Zuweisung auf einen Kreis, die die Behandlung
	%definiert. Der im Erhebungsjahr gemessene Status ist damit f\"ur genau die
	%wechselnden Personen ein Nachbehandlungsmerkmal. Die Sch\"atzung wird
	%deshalb zweifach berichtet -- mit Status im Erhebungsjahr (vergleichbar zur
	%Vorlage) und mit Status bei Erstbeobachtung (konzeptionell sauberer); siehe
	%Tabelle~\ref{tab:ext-asyl-erwerb}.
}
\end{table}
